
\subsection{Busqueda binaria con decimales}

\begin{frame}[plain]
	\centering
	\color{theme-brand-color}
	\Huge \textbf{Busqueda binaria con decimales}
\end{frame}

\begin{frame}{Busqueda binaria con decimales}
	Si tenemos \(l = 0.0\) y \(r = 10.0\), donde \(l, r \in \mathbb{R}\), habrá un número infinito de
	numeros reales en el intervalo \([l, r]\). Pero en números decimales con dos cifras después del punto hay...
\end{frame}

\begin{frame}{Busqueda binaria con decimales}
	En este caso tenemos una presición de $10^{-2}$, entonces si bien hay 10 enteros entre $1$ y $10$, si contamos la presición tendremos un total de $1000$ números:

	\begin{align*}
		1.00   \\
		1.01   \\
		\vdots \\
		2.00   \\
		2.01   \\
		\vdots \\
		10.00
	\end{align*}
	En general, podemos decir que si buscamos presición de $\epsilon$, entonces hace falta considerar $\frac{r - l}{\epsilon}$ números
\end{frame}

\begin{frame}{Búsqueda binaria con decimales}
	Comúnmente hay problemas donde nos dicen que la respuesta será considerada si el error absoluto no es más grande que $10^{-6}$. Entonces, si por ejemplo tenemos:
	\[
		\begin{aligned}
			l        & = 0.0     \\
			r        & = 10^{9}  \\
			\epsilon & = 10^{-6}
		\end{aligned}
	\]

	Consideraremos un total de:
	\[
		\dfrac{r - l}{\epsilon} = \dfrac{10^{9} - 0}{10^{-6}} = 10^{9} \cdot 10^{6} = 10^{15}
	\]
\end{frame}

\begin{frame}{¿Cómo hacemos binary search con decimales?}
	Sabemos que el algoritmo de binary search divide el espacio de busqueda en $2$ cada vez, y con lo anterior vimos que nuestro espacio de busqueda es de $\frac{r-l}{\epsilon}$ números, entonces binary search tomará $\log_2(\frac{r-l}{\epsilon})$ pasos en llegar a la respuesta.

	\vspace{0.5cm}

	Entonces en vez de hacer un \texttt{while($l \leq r$){$\dots$}}, es mejor hacer un número de iteraciones correspondinete al valor de \(\log_2({\frac{r-l}{\epsilon}})\)
\end{frame}


\begin{frame}{Problema ejemplo}
	\begin{ExampleP}[title=Equation]
		Encuentra el número $x$ tal que $x^{2} + \sqrt{x} = c$.

		\textbf{Input:}

		Solo un número $c \in \mathbb{R}$ tal que $1.0 \leq c \leq 10^{10}$

		\textbf{Output: }

		Imprime el número $x$ requerido. La respuesta será considerada correcta si el error absoluto no es más grande que $10^{-6}$
	\end{ExampleP}

	\begin{ExampleP}
		\textbf{input: }
		$15.6$

		\textbf{output: }
		$3.698232168829691$
	\end{ExampleP}
\end{frame}

\begin{frame}[fragile]{Solución a 'Equation'}
	\begin{center}
		\begin{minipage}{0.8\textwidth}
			\inputminted[
				firstline=1,
				lastline=24,
				linenos,
				fontsize=\tiny,
				numbersep=3pt,
				breaklines
			]{cpp}{code/binary_search_on_the_answer/equation.cpp}
		\end{minipage}
	\end{center}
\end{frame}

\begin{frame}

	\begin{TipP}

		\textbf{Observación: } Se hace $l = mid$ y $r = mid$ y no $l = mid \pm 1$ o $r = mid \pm 1$ pues eso nos alejaría más de la respuesta
	\end{TipP}

	\begin{TipP}
		\textbf{Nota: } En general, con un número de iteraciones de $60$ para el \texttt{for} es suficiente para muchos problemas, pero se acostumbra a poner $100$ o $200$
	\end{TipP}
\end{frame}
